% SolarMax - Pre-Deployment Checklist (point it at the sun)
% Compile:  pdflatex deployment_checklist.tex   (run twice for the table of contents)

\documentclass[11pt]{article}
\usepackage[margin=2.2cm]{geometry}
\usepackage{xcolor}
\usepackage[most]{tcolorbox}
\usepackage{enumitem}
\usepackage{titlesec}
\usepackage{listings}
\usepackage{amssymb}   % \square for the checklist boxes
\usepackage[colorlinks=true,linkcolor=blue!60!black,urlcolor=blue!60!black]{hyperref}

% ── Colors / section styling ─────────────────────────────────────────────────
\definecolor{brandblue}{HTML}{1565C0}
\definecolor{brandgreen}{HTML}{2E7D32}
\titleformat{\section}{\Large\bfseries\color{brandblue}}{\thesection.}{0.5em}{}

% ── Callout boxes ─────────────────────────────────────────────────────────────
\newtcolorbox{tip}{colback=green!6,colframe=brandgreen,boxrule=0.8pt,
  left=8pt,right=8pt,top=6pt,bottom=6pt,title=\textbf{Tip},fonttitle=\color{white}}
\newtcolorbox{warn}{colback=red!5,colframe=red!70!black,boxrule=0.8pt,
  left=8pt,right=8pt,top=6pt,bottom=6pt,title=\textbf{Important},fonttitle=\color{white}}
\newtcolorbox{note}{colback=blue!5,colframe=brandblue,boxrule=0.8pt,
  left=8pt,right=8pt,top=6pt,bottom=6pt,title=\textbf{Note},fonttitle=\color{white}}

% ── Code / command / serial-output blocks ────────────────────────────────────
\lstdefinestyle{cmd}{basicstyle=\ttfamily\small,backgroundcolor=\color{black!6},
  frame=single,rulecolor=\color{black!25},framesep=6pt,aboveskip=6pt,belowskip=6pt,
  xleftmargin=4pt,xrightmargin=4pt,breaklines=true,columns=fullflexible,keepspaces=true}
\lstset{style=cmd}

\setlist[enumerate]{leftmargin=1.7em,itemsep=3pt,topsep=4pt}

\begin{document}

% ── Title ─────────────────────────────────────────────────────────────────────
\begin{center}
{\Huge\bfseries\color{brandblue} SolarMax}\\[4pt]
{\LARGE Pre-Deployment Checklist --- Point It at the Sun}\\[8pt]
{\large Four numbers to check before the panel goes outside.}\\[6pt]
\rule{0.9\linewidth}{0.8pt}
\end{center}

\vspace{4pt}
\tableofcontents
\vspace{8pt}

% =============================================================================
\section{Before you start}
This assumes you've already \textbf{flashed} the board and \textbf{set the clock}
(see the Clock Setup guide). This guide covers the last step: making sure the panel
actually points at the \emph{right} place when you power it up outside.

\begin{note}
\textbf{How the tracking works:} SolarMax computes where the sun is purely from the
\textbf{date, time, and location} --- there are no light sensors steering it. So on
power-up it works out the sun's position and drives the panel there automatically.
That means it can only be right if the \textbf{location} and the \textbf{mounting
angle} are correct (both in \texttt{include/config.h}), and the \textbf{limit
switches} are working.
\end{note}

\begin{note}
\textbf{Position is tracked by the motor and limit switches --- no calibration step
for you.} On power-up the panel homes to the limit switches and measures its own
travel speed automatically (you'll see a \texttt{[MOTOR] Calibrated:\ \ldots\ deg/s}
line). The potentiometer is \textbf{not} used for this anymore --- it stays wired but
does nothing for control.
\end{note}

% =============================================================================
\section{First --- get the updated code onto the computer}
The firmware was \textbf{updated} so the panel tracks its position with the motor and
limit switches instead of the potentiometer --- that's the fix for the repeating
\texttt{move timed out --- check pot wiring} errors. Get that new code onto your
computer before flashing:

\begin{itemize}[leftmargin=1.6em]
  \item \textbf{If you cloned it with Git:} open the PlatformIO terminal and run
        \texttt{git pull} inside the project folder.
  \item \textbf{If you downloaded a ZIP} (the flashing-guide way): download the latest
        ZIP from the repo again and replace your project folder with it.
\end{itemize}

\begin{warn}
Keep your own \texttt{config.h} settings --- your Wi-Fi (lines 9--10) and location
(lines 4--5). If you replaced the whole folder, re-apply those edits before flashing.
\end{warn}

\begin{note}
Re-flashing the \emph{old} code won't help --- it still expects the pot. You're on the
updated code when it prints \texttt{[MOTOR] Homing to east limit\ldots} on boot.
\end{note}

% =============================================================================
\section{Check 1 --- Location (config.h lines 4--5)}
Open \texttt{include/config.h} and set your site's coordinates on
\textbf{lines 4 and 5}:
\begin{lstlisting}
#define LATITUDE          33.4255f    // Decimal degrees North
#define LONGITUDE       -111.9400f    // Decimal degrees (negative = West)
\end{lstlisting}
The defaults are Tempe, AZ. Change them to wherever the panel will actually sit.

\begin{tip}
\textbf{Getting your coordinates:} in Google Maps, right-click your exact spot ---
the latitude and longitude appear at the top of the menu (e.g. \texttt{33.4255,
-111.9400}); click them to copy. \textbf{North latitude is positive}; \textbf{West
longitude is negative} (keep the minus sign). Keep the trailing \texttt{f}.
\end{tip}

% =============================================================================
\section{Check 2 --- Mounting angle (config.h lines 86--87)}
These describe how the rotation \emph{axis} is mounted, on \textbf{lines 86 and 87}:
\begin{lstlisting}
#define AXIS_TILT_DEG      0.0f   // Degrees above horizontal (0 = flat roof)
#define AXIS_AZIMUTH_DEG   0.0f   // 0=North, 90=East, 180=South, 270=West
\end{lstlisting}

\begin{itemize}[leftmargin=1.6em]
  \item \textbf{Standard install (flat and level, axis running North--South):} leave
        both at \texttt{0.0f}. This is the common case --- if in doubt, use this.
  \item \textbf{If the mount is tilted} (sloped roof, or a seasonal tilt): set
        \texttt{AXIS\_TILT\_DEG} to the pitch angle (degrees above horizontal), and
        \texttt{AXIS\_AZIMUTH\_DEG} to the compass direction the \emph{raised} end
        of the axis faces.
\end{itemize}

\begin{note}
Getting these wrong tilts the whole day's tracking off. If you're not tilting the
mount, \texttt{0 / 0} is correct. The full geometry is in
\texttt{docs/sun\_tracking\_math.md} if you need it.
\end{note}

% =============================================================================
\section{Check 3 --- Limit switches (position calibrates itself)}
There is \textbf{no manual calibration} here. Position is tracked by the motor and
the two limit switches: on power-up the panel homes to the limits and measures its
own travel speed automatically. Your only job is to confirm the \textbf{limit
switches actually work}.

When you power up (Serial Monitor open) you should see, within a minute:
\begin{lstlisting}
[MOTOR] Homing to east limit...
[MOTOR] Sweeping to west limit to measure travel speed...
[MOTOR] Calibrated: 60 deg in 18.4 s = 3.26 deg/s
\end{lstlisting}
Seeing the \texttt{Calibrated} line means both limit switches tripped and the panel
knows where it is. You're done --- nothing to edit.

\begin{warn}
If instead you see \texttt{never reached east/west limit --- check limit-switch
wiring}, a limit switch isn't triggering. The tracker will \textbf{not} move until
this is fixed --- the limit switches are what tell it where the panel is. Check their
wiring (and \texttt{LIMIT\_ACTIVE} in \texttt{config.h} if you use normally-closed
switches).
\end{warn}

\begin{note}
The panel does a slow home-and-sweep every time it boots. In the field it never
powers off, so you'd only ever notice this while (re)programming.
\end{note}

% =============================================================================
\section{Finish --- upload and power up outside}
\begin{enumerate}
  \item Re-flash the real firmware after your edits:
\begin{lstlisting}
pio run -e esp32dev -t upload
\end{lstlisting}
  \item Power the panel up \textbf{during daylight}. Within a few seconds it computes
        the sun's position and drives the panel to face it, then nudges it westward
        through the day. In the Serial Monitor you'll see lines like:
\begin{lstlisting}
[SUN]  Local 10:15  El=42.7 deg  Az=131.6 deg  Panel target=-8.4 deg
\end{lstlisting}
\end{enumerate}

\begin{tip}
\textbf{Plugged in at night?} That's fine --- it parks and prints
\texttt{[NIGHT] Waiting for sunrise\ldots}, then starts facing the sun automatically
at sunrise (sun above \texttt{SUN\_MIN\_ELEV\_DEG}, config.h line 94). High wind
(\(\geq\) \texttt{WIND\_STOW\_MPH}, config.h line 92) sends it flat to stow.
\end{tip}

% =============================================================================
\section{Quick checklist}
\begin{itemize}[leftmargin=2.2em,itemsep=5pt]
  \item[$\square$] Clock set (Clock Setup guide) --- no \texttt{No valid time source} error.
  \item[$\square$] \textbf{Latitude} correct --- \texttt{config.h} line 4.
  \item[$\square$] \textbf{Longitude} correct (minus sign if West) --- \texttt{config.h} line 5.
  \item[$\square$] \textbf{Axis tilt} matches the mount (\texttt{0} if flat) --- \texttt{config.h} line 86.
  \item[$\square$] \textbf{Axis azimuth} matches the mount (\texttt{0} if flat) --- \texttt{config.h} line 87.
  \item[$\square$] \textbf{Limit switches work} --- saw \texttt{[MOTOR] Calibrated:\ \ldots\ deg/s} on boot (Check~3).
  \item[$\square$] Re-flashed \texttt{esp32dev}, powered up in daylight, panel faces the sun.
\end{itemize}

% =============================================================================
\section{Troubleshooting --- ``it moves, but to the wrong place''}
\begin{itemize}[leftmargin=1.6em]
  \item \textbf{Doesn't move at all, logs \texttt{never reached\ldots limit} or
        \texttt{not calibrated}:} a limit switch isn't triggering, so the panel never
        homed --- check the limit-switch wiring (Check~3).
  \item \textbf{Panel drifts a few degrees off during the day:} normal for open-loop
        timing --- it re-zeroes itself at the limit switches each morning and evening.
  \item \textbf{Tracking is off by a consistent amount all day:} the mounting angle
        is wrong --- recheck Check~2 (config.h lines 86--87), and that the axis really
        runs North--South.
  \item \textbf{Off by hours / points east when it should point west:} location is
        wrong --- recheck Check~1 (config.h lines 4--5), especially the longitude sign.
  \item \textbf{Doesn't move at all, logs \texttt{No valid time source}:} the clock
        isn't set --- see the Clock Setup guide.
  \item \textbf{Sits flat and won't track:} it's in stow (wind) or night --- check the
        wind reading and that the sun is actually up.
\end{itemize}

\end{document}
